\section{Methodology}\label{sec:methodology}

This section presents the air-cut referencing methodology for extracting toolpath fingerprints from CNC sensor data. We describe the data acquisition pipeline, temporal binning strategy, signal isolation methods, dimensionality reduction, and classification framework.

\subsection{Data Acquisition and Sensor Configuration}\label{sec:data_acquisition}

We instrument a desktop CNC milling machine with two categories of sensors: six wireless IMU (inertial measurement unit) modules mounted at distinct locations on the machine frame and moving components, and four electrical sensors measuring spindle and axis motor currents. Each IMU module records 17 channels comprising three-axis linear acceleration, three-axis angular velocity, three-axis magnetometer readings, and derived orientation quantities. The electrical sensors capture spindle motor current and current draw plus acceleration for the X, Y, and Z linear axes. Together, these provide 110 features per timestep across all sensor modules.

Data is collected at an approximately \SI{0.25}{\second} sampling interval, synchronized across all sensor modules via timestamp alignment. Each experimental run produces a single CSV file containing the time-aligned multi-sensor data stream for the full duration of the toolpath execution.

We collect data under two conditions:
\begin{itemize}
    \item \textbf{Air-cut runs}: The CNC machine executes the toolpath program with the spindle running but with no workpiece mounted, so the tool moves through air without engaging any material.
    \item \textbf{Active-cutting runs}: The machine executes the toolpath on a TIVAR UHMW-PE (ultra-high molecular weight polyethylene) workpiece, producing actual material removal.
\end{itemize}

Three toolpath types are studied: \emph{adaptive clearing}, \emph{facing}, and \emph{pocketing}. These represent fundamentally different engagement geometries and cutting force profiles, providing a meaningful classification challenge.

\subsection{Temporal Binning and Feature Extraction}\label{sec:temporal_binning}

\subsubsection{Motivation: Why G-Code Alignment Fails}

An ideal approach to air-cut referencing would align the air-cut and active-cutting signals on a per-G-code-command basis, enabling direct subtraction at each point in the toolpath. However, this requires that both datasets be generated from the same G-code program with identical command sequences. In our experimental setting, the air-cut and active-cutting programs were generated by different CAM configurations, resulting in different G-code command sequences for the same nominal toolpath type. Feed rates, plunge strategies, and retract moves differ between programs, making point-by-point temporal alignment infeasible.

\subsubsection{Temporal Binning Strategy}

We address this misalignment through temporal binning, which captures the meso-scale temporal structure of each run without requiring command-level synchronization. Given a sensor data matrix $\mathbf{X} \in \mathbb{R}^{T \times C}$ for a run of $T$ timesteps and $C$ channels, we divide the time axis into $B$ equal-duration bins:
\begin{equation}
    \mathbf{X}^{(b)} = \left\{ \mathbf{x}_t \in \mathbb{R}^C \;\middle|\; t \in \left[\frac{(b-1) \cdot T}{B},\; \frac{b \cdot T}{B}\right) \right\}, \quad b = 1, \ldots, B
\end{equation}
where $\mathbf{X}^{(b)}$ denotes the set of timesteps falling within bin~$b$.

\subsubsection{Statistical Feature Computation}

For each bin $b$ and each channel $c$, we compute seven summary statistics:
\begin{equation}
    \mathbf{f}^{(b,c)} = \begin{bmatrix}
        \mu_{b,c} \\[2pt] \sigma_{b,c} \\[2pt] \min_{b,c} \\[2pt] \max_{b,c} \\[2pt] \text{median}_{b,c} \\[2pt] \text{skew}_{b,c} \\[2pt] \text{kurt}_{b,c}
    \end{bmatrix}
    = \begin{bmatrix}
        \frac{1}{|\mathbf{X}^{(b)}|} \sum_{t \in b} x_{t,c} \\[4pt]
        \sqrt{\frac{1}{|\mathbf{X}^{(b)}|-1} \sum_{t \in b} (x_{t,c} - \mu_{b,c})^2} \\[4pt]
        \min_{t \in b} x_{t,c} \\[4pt]
        \max_{t \in b} x_{t,c} \\[4pt]
        \text{median}(\{x_{t,c}\}_{t \in b}) \\[4pt]
        \frac{1}{|\mathbf{X}^{(b)}|} \sum_{t \in b} \left(\frac{x_{t,c} - \mu_{b,c}}{\sigma_{b,c}}\right)^3 \\[4pt]
        \frac{1}{|\mathbf{X}^{(b)}|} \sum_{t \in b} \left(\frac{x_{t,c} - \mu_{b,c}}{\sigma_{b,c}}\right)^4 - 3
    \end{bmatrix}
\end{equation}
where $|\mathbf{X}^{(b)}|$ is the number of timesteps in bin~$b$.

\subsubsection{Run-Level Feature Vector Assembly}

Each of the $7C$ per-bin statistics forms a length-$B$ trajectory across the temporal bins. To obtain a fixed-length run-level descriptor whose dimensionality does not grow with the bin count, we summarize each per-bin statistic by its \emph{mean} and \emph{standard deviation across the $B$ bins}:
\begin{equation}\label{eq:feature_vector}
    \mathbf{v} = \Big[\, \operatorname{mean}_{b}\big(f^{(b,c)}_s\big),\;\; \operatorname{std}_{b}\big(f^{(b,c)}_s\big) \,\Big]_{\,c=1,\ldots,C;\; s=1,\ldots,7} \in \mathbb{R}^{2 \cdot 7 \cdot C},
\end{equation}
where $f^{(b,c)}_s$ is the $s$-th of the seven statistics for bin~$b$ and channel~$c$. The mean-across-bins captures each statistic's average level over the program, while the std-across-bins captures its temporal variation. This yields a run-level feature vector of dimensionality $d = 14C$; for $C = 110$ channels, $d = 1{,}540$.

The whole-run baseline corresponds to $B = 1$: the seven statistics are computed once over all timesteps, with no across-bin variation, yielding $d = 7C = 770$. We use this to isolate the benefit of temporal structure---the additional $770$ std-across-bins features that encode how each statistic evolves through the program.

\subsection{Air-Cut Referencing and Signal Isolation}\label{sec:signal_isolation}

Given the temporal-binned feature vectors for active-cutting runs and the corresponding air-cut baseline statistics, we define four signal isolation methods of increasing sophistication.

Let $\mathbf{v}_\text{active} \in \mathbb{R}^d$ denote the feature vector of an active-cutting run and let $\boldsymbol{\mu}_\text{air} \in \mathbb{R}^d$ and $\boldsymbol{\sigma}_\text{air} \in \mathbb{R}^d$ denote the element-wise mean and standard deviation of the air-cut feature vectors across all air-cut runs, computed independently for each feature dimension.

\subsubsection{Method 1: Raw (No Referencing)}

The raw method uses the active-cutting feature vector without any air-cut referencing:
\begin{equation}\label{eq:raw}
    \mathbf{v}_\text{raw} = \mathbf{v}_\text{active}
\end{equation}
This serves as the baseline, representing conventional feature extraction without signal isolation.

\subsubsection{Method 2: Mean-Subtracted}

The subtracted method removes the average air-cut contribution from each feature:
\begin{equation}\label{eq:subtracted}
    \mathbf{v}_\text{sub} = \mathbf{v}_\text{active} - \boldsymbol{\mu}_\text{air}
\end{equation}
This centers the features relative to the air-cut baseline, removing the mean machine-specific offset. Features that are identical during air cutting and active cutting will be driven toward zero, while features that differ---those arising from the cutting process---will retain nonzero values.

\subsubsection{Method 3: Ratio-Normalized}

The ratio method normalizes each feature by the corresponding air-cut mean:
\begin{equation}\label{eq:ratio}
    \mathbf{v}_\text{ratio} = \frac{\mathbf{v}_\text{active}}{\boldsymbol{\mu}_\text{air}}
\end{equation}
where division is element-wise. Features that are unchanged by cutting will have ratio values near unity, while features amplified or attenuated by the cutting process will deviate proportionally. This method is scale-invariant: sensors with different absolute magnitudes are placed on a common footing. Features where $\mu_{\text{air},j} \approx 0$ are excluded or clamped to avoid division instability.

\subsubsection{Method 4: Z-Score-Normalized}

The z-score method normalizes each feature by both the mean and standard deviation of the air-cut distribution:
\begin{equation}\label{eq:zscore}
    \mathbf{v}_\text{z} = \frac{\mathbf{v}_\text{active} - \boldsymbol{\mu}_\text{air}}{\boldsymbol{\sigma}_\text{air}}
\end{equation}
This expresses each feature as the number of air-cut standard deviations by which the active-cutting value differs from the air-cut mean. Features with high natural variability during air cutting (e.g., due to measurement noise or spindle vibration harmonics) are appropriately down-weighted, while features that are stable during air cutting but change substantially during active cutting receive high z-scores. Features where $\sigma_{\text{air},j} < \epsilon$ (for a small threshold $\epsilon$) are excluded to avoid amplification of near-constant channels.

\subsection{Dimensionality Reduction}\label{sec:dim_reduction}

The high dimensionality of the run-level feature vectors ($d = 1{,}540$ for 110~channels, a $\sim$30:1 feature-to-sample ratio) relative to the number of available runs ($n = 51$ active-cutting runs) creates a risk of overfitting. We employ two complementary dimensionality reduction strategies:

\textbf{Principal Component Analysis (PCA).} We apply PCA to the feature matrix, retaining the minimum number of components that explain at least 95\% of the total variance. PCA is applied within the cross-validation loop to prevent information leakage from the test folds.

\textbf{Random Forest Feature Importance.} We use the Gini importance scores from a preliminary Random Forest fit to rank features and select the top $k$ features. This supervised approach can identify features that are specifically discriminative for the classification task, complementing the unsupervised variance-based selection of PCA.

Both methods are evaluated independently to assess whether supervised or unsupervised dimensionality reduction is more effective for this task.

\subsection{Classification Framework}\label{sec:classification}

We evaluate three classifiers representing different modeling paradigms:

\textbf{Random Forest (RF).} An ensemble of 500 decision trees with Gini impurity splitting, bootstrap aggregation, balanced class weights, minimum 2 samples per leaf, and random state 42~\cite{breiman2001randomforests}.

\textbf{Support Vector Machine (SVM).} A multi-class SVM with RBF kernel~\cite{widodo2007svm}, $C = 1.0$, $\gamma = \text{scale}$ (i.e., $\gamma = 1 / (n_\text{features} \cdot \mathrm{Var}(X))$), and balanced class weights.

\textbf{Multi-Layer Perceptron (MLP).} Three hidden layers of 256, 128, and 64 neurons with ReLU activations, L2 regularization ($\alpha = 0.01$), Adam optimizer (default learning rate $10^{-3}$), maximum 2{,}000 epochs with early stopping (patience on 15\% validation split), and random state 42.

The z-score isolation method uses a safety threshold $\epsilon = 10^{-6}$: when the air-cut standard deviation $\sigma_\text{air} < \epsilon$ for a given feature-bin combination, the z-score is set to zero rather than producing a division-by-zero artifact.

\subsubsection{Cross-Validation Protocol}

All classification experiments use stratified $k$-fold cross-validation with $k = 5$. Critically, the stratification is performed at the \emph{run level}, ensuring that all timesteps and features from a single run appear entirely in either the training or test fold, never split across folds. This prevents temporal leakage that would occur if consecutive timesteps from the same run appeared in both training and test sets.

Within each fold, the training data is used to compute the air-cut baseline statistics ($\boldsymbol{\mu}_\text{air}$, $\boldsymbol{\sigma}_\text{air}$), fit the PCA transformation, and train the classifier. The test fold's active-cutting runs are then transformed using the training fold's air-cut statistics and PCA projection before prediction. This ensures that the entire pipeline---from air-cut referencing through classification---is evaluated without data leakage.

Performance is reported as the mean and standard deviation of accuracy, precision, recall, and F1 score across folds, along with the aggregate confusion matrix.
